\section{Aplicação das avaliações}\label{sec:aplicacoes}

A aplicação relaciona um instrumento avaliativo a um aprendente e registra as respostas produzidas durante sua utilização. A estrutura de dados correspondente (entidades \textit{aplicacoes} e \textit{respostas}) foi descrita na Seção~\ref{sec:modelagem}; esta seção descreve o andamento de uma aplicação.

Para representar esse andamento, foram definidos quatro estados com transições bem delimitadas: \textit{PENDING} (criada, sem respostas registradas), \textit{IN\_PROGRESS} (respostas salvas em modo de progresso normal), \textit{DRAFT} (respostas salvas explicitamente como rascunho pelo profissional, sinalizando preenchimento pausado intencionalmente) e \textit{COMPLETED} (finalizada). A distinção entre \textit{IN\_PROGRESS} e \textit{DRAFT} é determinada por uma escolha explícita do profissional no momento de salvar, o que possibilita interromper o preenchimento sem perder os dados já registrados e retomá-lo posteriormente.

A partir de qualquer estado não finalizado, o profissional pode concluir a aplicação. Nesse momento, o sistema verifica o preenchimento das questões obrigatórias, e as respostas ficam bloqueadas para edição; aplicações finalizadas também não podem ser excluídas. O formulário de preenchimento é construído a partir da configuração das questões do instrumento, apresentando para cada uma o componente de entrada correspondente ao seu tipo.

A separação entre o cadastro do instrumento e suas aplicações permite preservar o histórico dos procedimentos realizados: o instrumento funciona como uma estrutura reutilizável, enquanto cada aplicação representa uma ocorrência específica de sua utilização, associada à versão do instrumento empregada (Seção~\ref{sec:instrumentos}). Essa organização é importante para o acompanhamento longitudinal do aprendente, pois diferentes aplicações podem ser associadas ao mesmo indivíduo em momentos distintos, mantendo seus respectivos registros.
